% ৪.১০ ফর্ম ও ডেটা সংগ্রহ
\section{ফর্ম ও ডেটা সংগ্রহ}

\begin{learningbox}
এই টপিক শেষে শিক্ষার্থী পারবে:
\begin{itemize}
\item HTML form-এর কাজ ও মৌলিক tag লিখতে।
\item input type, textarea, select, button ব্যবহার করতে।
\item GET ও POST method-এর পার্থক্য ব্যাখ্যা করতে।
\end{itemize}
\end{learningbox}

\subsection{Form}
ব্যবহারকারীর কাছ থেকে data নেওয়ার জন্য HTML form ব্যবহৃত হয়। registration, login, search, feedback, admission application ইত্যাদি form-এর উদাহরণ।

\begin{definitionbox}{Form}
যে HTML structure ব্যবহারকারীর input গ্রহণ করে server বা script-এর কাছে পাঠাতে পারে, তাকে form বলা হয়।
\end{definitionbox}

\begin{center}
\begin{tabular}{|p{0.22\textwidth}|p{0.42\textwidth}|p{0.22\textwidth}|}
\hline
\textbf{Tag/Attribute} & \textbf{কাজ} & \textbf{উদাহরণ} \\
\hline
\texttt{<form>} & form container & action, method থাকে \\
\hline
\texttt{<input>} & এক লাইনের input & text, password, email \\
\hline
\texttt{<textarea>} & বহু লাইনের text & comment/message \\
\hline
\texttt{<select>} & dropdown list & option list \\
\hline
\texttt{required} & input বাধ্যতামূলক করে & validation \\
\hline
\end{tabular}
\end{center}

\begin{center}
\fbox{\parbox{0.9\textwidth}{
\textbf{Form example placeholder}\\[4pt]
\texttt{<form action="submit.php" method="post">}\\
\texttt{<input type="text" name="name" required>}\\
\texttt{<input type="email" name="email">}\\
\texttt{<button type="submit">Send</button>}\\
\texttt{</form>}
}}
\end{center}

\subsection{GET ও POST}
\begin{center}
\begin{tabular}{|p{0.25\textwidth}|p{0.30\textwidth}|p{0.30\textwidth}|}
\hline
\textbf{ভিত্তি} & \textbf{GET} & \textbf{POST} \\
\hline
Data পাঠানো & URL-এর সাথে query হিসেবে & request body-তে \\
\hline
ব্যবহার & search, filter & login, registration, sensitive data \\
\hline
গোপনীয়তা & কম & তুলনামূলক ভালো, তবে HTTPS দরকার \\
\hline
\end{tabular}
\end{center}

\begin{mistakebox}{সাধারণ ভুল}
Client-side validation সুবিধাজনক, কিন্তু নিরাপত্তার জন্য server-side validation দরকার। শুধু browser validation-এর উপর নির্ভর করা ঠিক নয়।
\end{mistakebox}
